\chapter{Calibration}
\label{chap:calibration}

\section{Calibration Architecture}
\label{sec:calib-architecture}

% TODO: The build/apply split -- calibration/ builds artifacts,
% preprocessing/ applies them -- and why that separation exists (a science
% frame's settings are checked against the artifact's recorded
% exposure/gain/timestamp via CalibrationRecord before it's applied). Frame
% this as the organizing principle for the rest of the chapter.

\section{Sensor Calibration}
\label{sec:calib-sensor}

\subsection{Baseline}
\label{sec:calib-baseline}

% TODO: Per-session background/dark average; background_sigma measured for
% free from the same stacked frames, used later in the noise model
% (chapter 4).

\subsection{Flat Field}
\label{sec:calib-flat-field}

% TODO: PRNU correction from uniform-illumination frames, dark-subtracted
% first; rejects a saturated source outright.

\subsection{Bad Pixel Map}
\label{sec:calib-bad-pixel}

% TODO: Dead/hot pixels flagged from flat-field outliers; zeroed rather
% than interpolated when applied (exact for a weighted centroid).

\subsection{Conversion Gain}
\label{sec:calib-conversion-gain}

% TODO: Photon transfer curve -- uniform illumination at fixed brightness,
% swept across exposure times; variance computed temporally (not
% spatially) so PRNU doesn't bias the gain low.

\section{Spatial Calibration}
\label{sec:calib-spatial}

% TODO: Fixed scale factor (ratio of the two relay-lens focal lengths)
% rather than a translation-stage fit, and why that's precise enough for
% this project's scope. Practical measured magnification M = 1.504 +/-
% 0.008 -- see presentation slide 24.

\section{Spectral Calibration}
\label{sec:calib-spectral}

\subsection{Grating Geometry and Predicted Line Spacing}
\label{sec:calib-grating-geometry}

% TODO: Transmission-grating equation predicting relative pixel spacing
% between known Argon lines, used to search for/identify peaks.

\subsection{Line Matching}
\label{sec:calib-line-matching}

% TODO: Peak detection (1D collapse + find_peaks + sub-pixel refinement)
% and matching detected peaks to reference Argon lines.

\subsection{Wavelength Fit and Non-Linear Dispersion}
\label{sec:calib-wavelength-fit}

% TODO: Linear/quadratic/cubic pixel-to-wavelength fits and the evidence
% for non-linear spectral dispersion (reduced chi-squared comparison,
% pixel-column/uncertainty correlation r = -0.96). See presentation
% slide 23.
